% =====================================================================
%  Solutions: Emergent Causal Order from Entanglement
%
%  Companion document to emergent_causal_order_homework.tex.  Each
%  problem is restated, then a full solution is given.  Section,
%  equation, and figure references in the body refer to the paper
%  (emergent_causal_order.tex) unless qualified.
%
%  Build: pdflatex emergent_causal_order_solutions.tex
% =====================================================================

\documentclass[11pt,a4paper]{article}

\usepackage[utf8]{inputenc}
\usepackage[T1]{fontenc}
\usepackage{lmodern}
\usepackage{microtype}
\usepackage[a4paper,margin=1in]{geometry}

\usepackage{amsmath,amssymb,amsthm,mathtools}
\usepackage{bm}
\usepackage{enumitem}
\usepackage{tikz}
\usetikzlibrary{arrows.meta,positioning,calc}
\usepackage{hyperref}
\hypersetup{
  colorlinks=true,
  linkcolor=blue!50!black,
  urlcolor=blue!50!black,
}

% Problem environment
\newcounter{problem}
\newenvironment{problem}[1][]{%
  \refstepcounter{problem}%
  \par\medskip\noindent
  \textbf{Problem \theproblem.}\if\relax\detokenize{#1}\relax\else\ \textit{(#1)}\fi\quad\ignorespaces
}{\par\smallskip}

\newenvironment{solution}{%
  \par\noindent\textbf{Solution.}\quad\ignorespaces
}{\par\medskip}

% Macros (mirror the paper)
\newcommand{\ket}[1]{\lvert #1 \rangle}
\newcommand{\bra}[1]{\langle #1 \rvert}
\newcommand{\braket}[2]{\langle #1 \mid #2 \rangle}
\newcommand{\norm}[1]{\lVert #1 \rVert}
\newcommand{\abs}[1]{\lvert #1 \rvert}
\newcommand{\Tr}{\operatorname{Tr}}
\newcommand{\diag}{\operatorname{diag}}
\newcommand{\Hilb}{\mathcal{H}}
\newcommand{\real}{\mathbb{R}}
\newcommand{\complex}{\mathbb{C}}
\newcommand{\maj}{\mathrm{maj}}
\newcommand{\LR}{\mathrm{LR}}
\newcommand{\cs}{\mathrm{cs}}
\newcommand{\preceqmaj}{\preceq_{\maj}}
\newcommand{\preceqLR}{\preceq_{\LR}}
\newcommand{\preceqcs}{\preceq_{\cs}}

\title{Solutions Set\\
\large \emph{Emergent Causal Order from Entanglement}}
\author{}
\date{}

\begin{document}
\maketitle

\noindent
This is the solutions companion to the homework set associated with
the paper \emph{Emergent Causal Order from Entanglement: A
Pedagogical Introduction}.  Each problem is restated briefly, then a
complete solution is given.  Numerical values are exact where
possible and quoted to four significant figures otherwise.

\tableofcontents

% =====================================================================
\section{Partial orders and Kendall correlation}

\begin{problem}
$S = \{a, b, c, d\}$ with two partial orders defined by their strict
relations: $\preceq_1$ has $a \prec b, a \prec c, a \prec d, b \prec d, c \prec d$;
$\preceq_2$ has $a \prec b, a \prec c, b \prec c, b \prec d, c \prec d$.
Draw the Hasse diagrams, classify all pairs, and compute $\tau_{12}$.
\end{problem}

\begin{solution}
The Hasse diagrams are:

\smallskip
\begin{tikzpicture}[node distance=10mm and 14mm,
  every node/.style={font=\small},
  edge/.style={-{Stealth[length=2mm]}, thick}]
  \node (a1) at (0, 0)    {$a$};
  \node (b1) at (-1.2, 1) {$b$};
  \node (c1) at (1.2, 1)  {$c$};
  \node (d1) at (0, 2)    {$d$};
  \draw[edge] (a1) -- (b1);
  \draw[edge] (a1) -- (c1);
  \draw[edge] (b1) -- (d1);
  \draw[edge] (c1) -- (d1);
  \node[below=4mm of a1] {$\preceq_1$};
  \node (a2) at (5, 0)   {$a$};
  \node (b2) at (5, 1)   {$b$};
  \node (c2) at (5, 2)   {$c$};
  \node (d2) at (5, 3)   {$d$};
  \draw[edge] (a2) -- (b2);
  \draw[edge] (b2) -- (c2);
  \draw[edge] (c2) -- (d2);
  \node[below=4mm of a2] {$\preceq_2$};
\end{tikzpicture}

For $\preceq_2$, $b \prec d$ given is implied by $b \prec c \prec d$;
the Hasse diagram only shows cover edges.

The six unordered pairs:
\begin{center}
\renewcommand{\arraystretch}{1.1}
\begin{tabular}{c|c|c|l}
pair & in $\preceq_1$ & in $\preceq_2$ & class \\\hline
$\{a, b\}$ & $a \prec_1 b$ & $a \prec_2 b$ & concordant \\
$\{a, c\}$ & $a \prec_1 c$ & $a \prec_2 c$ & concordant \\
$\{a, d\}$ & $a \prec_1 d$ & $a \prec_2 d$ & concordant \\
$\{b, c\}$ & incomparable  & $b \prec_2 c$ & $\preceq_2$ only \\
$\{b, d\}$ & $b \prec_1 d$ & $b \prec_2 d$ & concordant \\
$\{c, d\}$ & $c \prec_1 d$ & $c \prec_2 d$ & concordant \\
\end{tabular}
\end{center}

So $n_{cc} = 5$, $n_{dd} = 0$, $n_1^o = 0$, $n_2^o = 1$, and
$\tau_{12} = 5/(5 + 0) = 1$.  The two orders agree on every pair
they both relate; $\preceq_2$ is strictly richer.
\end{solution}

\begin{problem}
Show majorization is transitive on $\Delta_{d-1}$.
\end{problem}

\begin{solution}
Suppose $x \preceq y$ and $y \preceq z$.  By definition, for every
$k \in \{1, \ldots, d\}$,
\[
  \sum_{i=1}^{k} x_i^\downarrow \;\le\; \sum_{i=1}^{k} y_i^\downarrow,
  \qquad
  \sum_{i=1}^{k} y_i^\downarrow \;\le\; \sum_{i=1}^{k} z_i^\downarrow.
\]
Chaining the two inequalities at each $k$ gives
$\sum_{i=1}^{k} x_i^\downarrow \le \sum_{i=1}^{k} z_i^\downarrow$.
Since $x, y, z$ all sum to $1$, the inequality is automatically an
equality at $k = d$.  This is exactly the definition of $x \preceq z$.
\end{solution}

\begin{problem}
On $S = \{1, 2, 3\}$, construct partial orders with (a) $\tau = -1$,
(b) $\tau$ undefined, (c) $\tau = 0$.
\end{problem}

\begin{solution}\hfill
\begin{enumerate}[label=(\alph*),leftmargin=*]
\item $\tau = -1$.  Take $\preceq_a$ to be the total order
      $1 \prec 2 \prec 3$ and $\preceq_b$ to be the reversed total
      order $3 \prec 2 \prec 1$.  Each of the three pairs is
      comparable in both, but oppositely.  $n_{cc} = 0$, $n_{dd} = 3$,
      $\tau = -3/3 = -1$.
\item $\tau$ undefined.  Take $\preceq_a$ relating only
      $1 \prec_a 2$, and $\preceq_b$ relating only $1 \prec_b 3$.
      The pairs $\{1, 2\}$ and $\{1, 3\}$ are each comparable in
      exactly one order; $\{2, 3\}$ is incomparable in both.  The
      both-comparable subset is empty, so the Kendall formula of the
      paper (eq.~3) has the form $0/0$ and $\tau$ is not defined.
\item $\tau = 0$.  Take $\preceq_a$ relating $1 \prec 2$ and
      $1 \prec 3$ (the bottom is $1$, the other two are
      incomparable).  Take $\preceq_b$ relating $1 \prec 2$ and
      $3 \prec 1$.  Then transitivity of $\preceq_b$ gives
      $3 \prec_b 2$.  The three pairs:
      $\{1, 2\}$ is concordant (both $1 \prec 2$);
      $\{1, 3\}$ is discordant ($1 \prec_a 3$ but $3 \prec_b 1$);
      $\{2, 3\}$ is comparable in $\preceq_b$ only.
      So $n_{cc} = 1$, $n_{dd} = 1$, and $\tau = (1-1)/(1+1) = 0$.
\end{enumerate}
\end{solution}

% =====================================================================
\section{Quantum states and the Schmidt decomposition}

\begin{problem}
$\ket\psi = \tfrac{1}{\sqrt 3}(\ket{00} + \ket{01} + \ket{10})$.
Find $\Psi$, $\rho_A$, the Schmidt spectrum, and the entanglement
entropy of the cut $\{1\}\,|\,\{2\}$.
\end{problem}

\begin{solution}
Read off the coefficient matrix $\Psi_{ij}$ where $i$ labels the
first qubit and $j$ the second:
\[
  \Psi \;=\; \frac{1}{\sqrt 3}\begin{pmatrix} 1 & 1 \\ 1 & 0 \end{pmatrix}.
\]
Compute
\[
  \rho_A \;=\; \Psi \Psi^\dagger
       \;=\; \frac{1}{3}\begin{pmatrix} 1 & 1 \\ 1 & 0 \end{pmatrix}
                       \begin{pmatrix} 1 & 1 \\ 1 & 0 \end{pmatrix}
       \;=\; \frac{1}{3}\begin{pmatrix} 2 & 1 \\ 1 & 1 \end{pmatrix}.
\]
The trace is $1$ as it must be.  The determinant is $(2 - 1)/9 = 1/9$.
The eigenvalues solve $\lambda^2 - \lambda + 1/9 = 0$:
\[
  \lambda_\pm \;=\; \frac{1 \pm \sqrt{1 - 4/9}}{2}
                \;=\; \frac{3 \pm \sqrt 5}{6}
                \;\approx\; (0.873,\, 0.127).
\]
The Schmidt spectrum is therefore
$\bm\lambda = \bigl(\tfrac{3 + \sqrt 5}{6},\; \tfrac{3 - \sqrt 5}{6}\bigr)$.
The entanglement entropy is
\[
  S = -\lambda_+ \ln \lambda_+ - \lambda_- \ln \lambda_-
    \;\approx\; -0.873 \cdot (-0.136) - 0.127 \cdot (-2.067)
    \;\approx\; 0.381 \text{ nats},
\]
or equivalently $\approx 0.550$ bits.
\end{solution}

\begin{problem}
Show by SVD that $\rho_A = \Psi\Psi^\dagger$ and
$\rho_B = \Psi^\dagger \Psi$ have identical nonzero spectra.
\end{problem}

\begin{solution}
Write $\Psi = U S V^\dagger$ with $U, V$ having orthonormal columns
and $S = \diag(s_1, \ldots, s_r)$ diagonal with $s_i > 0$.  Then
\[
  \rho_A = \Psi\Psi^\dagger = U S V^\dagger V S^\dagger U^\dagger
         = U (S S^\dagger) U^\dagger,
  \qquad
  \rho_B = \Psi^\dagger \Psi = V S^\dagger U^\dagger U S V^\dagger
         = V (S^\dagger S) V^\dagger.
\]
Both $S S^\dagger$ and $S^\dagger S$ are diagonal with diagonal
entries $(s_1^2, \ldots, s_r^2, 0, \ldots, 0)$ (extended by zeros to
the appropriate ambient dimension).  Conjugation by an isometry
preserves the spectrum, so $\rho_A$ and $\rho_B$ both have spectrum
$(s_1^2, \ldots, s_r^2, 0, \ldots, 0)$.  In particular their nonzero
spectra are identical.
\end{solution}

\begin{problem}[$\star$]
For $\ket\Psi = \cos\theta \ket{000} + \sin\theta \ket{111}$, compute
the Schmidt spectrum across the cuts $\{1\}\,|\,\{2,3\}$ and
$\{1,2\}\,|\,\{3\}$.
\end{problem}

\begin{solution}
The state is already in (a tensor-product version of) Schmidt form
for both cuts:
\begin{align*}
  \ket\Psi &= \cos\theta\, \ket{0}_1 \otimes \ket{00}_{23}
          + \sin\theta\, \ket{1}_1 \otimes \ket{11}_{23},\\
  \ket\Psi &= \cos\theta\, \ket{00}_{12} \otimes \ket{0}_3
          + \sin\theta\, \ket{11}_{12} \otimes \ket{1}_3.
\end{align*}
The vectors $\ket{0}_1, \ket{1}_1$ are orthonormal; so are $\ket{00}_{23}, \ket{11}_{23}$
and $\ket{00}_{12}, \ket{11}_{12}$ and $\ket{0}_3, \ket{1}_3$.  Each
expression is therefore already a Schmidt decomposition with
coefficients $(\cos\theta, \sin\theta)$.  The Schmidt spectrum is
$(\cos^2\theta, \sin^2\theta)$ in both cases.

Physically, the GHZ-like state is a single ``two-level'' bipartite
correlation that is shared identically across any non-trivial cut of
the chain.  No new Schmidt structure appears as more sites move to
one side of the cut.
\end{solution}

% =====================================================================
\section{Majorization}

\begin{problem}
Decide majorization comparability for four pairs.
\end{problem}

\begin{solution}\hfill
\begin{enumerate}[label=(\alph*),leftmargin=*]
\item Cumulative sums of $a = (0.5, 0.4, 0.1)$: $(0.5, 0.9, 1.0)$.
      Of $b = (0.5, 0.3, 0.2)$: $(0.5, 0.8, 1.0)$.  At every $k$,
      $a$'s cumulative sum is $\ge b$'s, so $b \preceq a$.
\item $a = (0.7, 0.2, 0.1)$ has cumsums $(0.7, 0.9, 1.0)$;
      $b = (0.6, 0.3, 0.1)$ has $(0.6, 0.9, 1.0)$.
      $a \succeq b$, i.e.\ $b \preceq a$.
\item $a = (0.4, 0.3, 0.2, 0.1)$ has cumsums $(0.4, 0.7, 0.9, 1.0)$;
      $b = (0.4, 0.4, 0.1, 0.1)$ has $(0.4, 0.8, 0.9, 1.0)$.
      Now $b$'s cumsum dominates, so $a \preceq b$.
\item $a = (0.5, 0.3, 0.2)$ has cumsums $(0.5, 0.8, 1.0)$;
      $b = (0.45, 0.40, 0.15)$ has $(0.45, 0.85, 1.0)$.  At $k = 1$,
      $a > b$.  At $k = 2$, $a < b$.  Neither cumulative sequence
      dominates the other.  $a$ and $b$ are \emph{incomparable}
      under majorization.
\end{enumerate}
\end{solution}

\begin{problem}
Why can the Bell pair not be deterministically converted into the
maximally entangled three-level state via LOCC?
\end{problem}

\begin{solution}
The Bell pair $\tfrac{1}{\sqrt 2}(\ket{00} + \ket{11})$ has Schmidt
spectrum $\bm\lambda_{\text{Bell}} = (\tfrac12, \tfrac12)$.
Zero-padded to length $3$ for comparison: $(\tfrac12, \tfrac12, 0)$.
The target state $\tfrac{1}{\sqrt 3}(\ket{00} + \ket{11} + \ket{22})$
has Schmidt spectrum $\bm\lambda_{\text{tgt}} = (\tfrac13, \tfrac13, \tfrac13)$.

Theorem~4.4 of the paper says LOCC convertibility
$\ket\psi \to \ket\phi$ requires
$\bm\lambda(\psi) \preceq \bm\lambda(\phi)$.  Compute cumulative
sums:
\[
  \bm\lambda_{\text{Bell}}: (\tfrac12,\, 1,\, 1),
  \qquad
  \bm\lambda_{\text{tgt}}: (\tfrac13,\, \tfrac23,\, 1).
\]
At $k = 1$, $\tfrac12 > \tfrac13$, so $\bm\lambda_{\text{Bell}}$ does
\emph{not} majorization-precede $\bm\lambda_{\text{tgt}}$.  Hence no
deterministic LOCC protocol exists.

Physically, the target state is more entangled (entropy $\ln 3$) than
the Bell pair (entropy $\ln 2$).  LOCC cannot increase entanglement,
and majorization is precisely the partial order that captures this
constraint.
\end{solution}

\begin{problem}[$\star$]
Show that majorization implies a reverse ordering on Shannon entropy:
$x \preceq y \Rightarrow H(x) \ge H(y)$.
\end{problem}

\begin{solution}
The function $f(p) = -p \ln p$ (with $f(0) := 0$) is concave on
$[0, 1]$: indeed $f''(p) = -1/p < 0$ on $(0, 1]$.  By the cited
Schur-concavity property, the symmetric function
$\Phi(x) = \sum_i f(x_i)$ inherits the property that
$x \preceq y \Rightarrow \Phi(x) \ge \Phi(y)$ for any concave $f$.
But $\Phi(x) = -\sum_i x_i \ln x_i = H(x)$.  Therefore
$H(x) \ge H(y)$ whenever $x \preceq y$.

In words: a more spread-out probability vector (lower in
majorization) has higher Shannon entropy.  Combined with Nielsen's
theorem, this gives the standard inequality that LOCC cannot
increase entanglement entropy.
\end{solution}

% =====================================================================
\section{The lattice Schwinger model}

\begin{problem}
Hopping operator $h_n = X_n X_{n+1} + Y_n Y_{n+1}$.  Write as a
$4\times 4$ matrix; show its action on the computational basis;
explain why it is called ``hopping''.
\end{problem}

\begin{solution}
With basis ordering $\{\ket{00}, \ket{01}, \ket{10}, \ket{11}\}$
(first index = first qubit), compute
\[
  X \otimes X = \begin{pmatrix} 0 & 0 & 0 & 1 \\ 0 & 0 & 1 & 0 \\ 0 & 1 & 0 & 0 \\ 1 & 0 & 0 & 0 \end{pmatrix},
  \qquad
  Y \otimes Y = \begin{pmatrix} 0 & 0 & 0 & -1 \\ 0 & 0 & 1 & 0 \\ 0 & 1 & 0 & 0 \\ -1 & 0 & 0 & 0 \end{pmatrix},
\]
using $Y\ket 0 = i \ket 1$ and $Y \ket 1 = -i\ket 0$ for the second
matrix.  Adding,
\[
  h \;=\; \begin{pmatrix} 0 & 0 & 0 & 0 \\ 0 & 0 & 2 & 0 \\ 0 & 2 & 0 & 0 \\ 0 & 0 & 0 & 0 \end{pmatrix}.
\]
The action on the basis is therefore
\[
  h\ket{00} = 0, \qquad h\ket{01} = 2\ket{10}, \qquad
  h\ket{10} = 2\ket{01}, \qquad h\ket{11} = 0.
\]
The operator swaps the two ``one-particle'' configurations
$\ket{01}$ and $\ket{10}$ (with a factor of $2$) and annihilates the
zero- and two-particle configurations.  Under the
Jordan--Wigner correspondence
$\ket 0 \leftrightarrow \text{empty},\, \ket 1 \leftrightarrow \text{occupied}$,
this is exactly the action of $c^\dagger_n c_{n+1} + c^\dagger_{n+1} c_n$
on the matter field: a single particle hops from one site to its
neighbor, with vacuum and doubly-occupied configurations annihilated
by the kinetic term.
\end{solution}

\begin{problem}
Compute $L_n$ for $n = 1, 2, 3$ on the all-$\ket 0$ configuration of
an $N = 4$ chain.
\end{problem}

\begin{solution}
With $L_0 = 0$ and $\langle Z_k \rangle = +1$ on every site,
$\tfrac{1 - Z_k}{2} = 0$ for all $k$, so
\[
  L_n \;=\; -\sum_{k=1}^{n} \frac{1 - (-1)^k}{2}.
\]
The summand is $1$ for odd $k$ and $0$ for even $k$, so $L_n$ counts
(with a sign) the number of odd $k$ in $\{1, \ldots, n\}$:
\[
  L_1 = -1,\qquad L_2 = -1,\qquad L_3 = -2.
\]
Physically, the all-$\ket 0$ configuration is \emph{not} the
half-filled staggered vacuum: every odd site is missing its
``positron'' (the c-number subtraction $\tfrac{1 - (-1)^k}{2}$
encodes the staggered vacuum reference).  The cumulative deficit
of $-1$ per odd site to the left of link $n$ is exactly $L_n$, the
net charge to the left of that link.  The corresponding electric
energy per link, $\tfrac{g^2 a}{2} L_n^2$, is non-zero, which is why
this configuration is not the ground state.
\end{solution}

\begin{problem}[$\star$]
Show $[M, H] = 0$ for $M = \sum_n Z_n$, where $H = H_{\mathrm{hop}} + H_m + H_E$.
\end{problem}

\begin{solution}
The pieces $H_m$ and $H_E$ are functions of the Pauli-$Z$ operators
only (the latter via the cumulative sums $L_n$).  Pauli-$Z$ operators
commute among themselves, so
$[Z_k, H_m] = 0$ and $[Z_k, H_E] = 0$ for every $k$, and therefore
$[M, H_m] = [M, H_E] = 0$.

For the hopping term the relevant identity is
$[Z_n + Z_{n+1},\, X_n X_{n+1} + Y_n Y_{n+1}] = 0$.  Using
$[Z_n, X_n] = 2iY_n$ and $[Z_n, Y_n] = -2i X_n$ (the basic
Pauli commutators) and $[A, BC] = [A, B]C + B[A, C]$:
\begin{align*}
  [Z_n,\, X_n X_{n+1}] &= 2iY_n X_{n+1}, &
  [Z_n,\, Y_n Y_{n+1}] &= -2i X_n Y_{n+1},\\
  [Z_{n+1},\, X_n X_{n+1}] &= 2i X_n Y_{n+1}, &
  [Z_{n+1},\, Y_n Y_{n+1}] &= -2i Y_n X_{n+1}.
\end{align*}
Adding the first column and the second column separately and then
combining,
\[
  [Z_n + Z_{n+1},\, X_n X_{n+1} + Y_n Y_{n+1}]
  = (2iY_n X_{n+1} - 2i X_n Y_{n+1}) + (2i X_n Y_{n+1} - 2i Y_n X_{n+1})
  = 0.
\]
Operators on disjoint sites commute, so $Z_k$ for any $k \notin \{n, n+1\}$
commutes with $X_n X_{n+1} + Y_n Y_{n+1}$ trivially.  Therefore
\[
  [M, H_{\mathrm{hop}}]
  \;=\; \frac{1}{4a}\sum_{n=1}^{N-1}
    [\sum_k Z_k,\, X_n X_{n+1} + Y_n Y_{n+1}]
  \;=\; 0.
\]
Combining: $[M, H] = 0$.  Translated to the fermion picture, this
says the total fermion number $\sum_n c^\dagger_n c_n = (N - M)/2$ is
conserved by the dynamics.
\end{solution}

% =====================================================================
\section{Lieb--Robinson and causal orders}

\begin{problem}
Five-site chain, snapshots at $t \in \{0, 0.5, 1.0, 1.5\}$.  Count
LR-related space-time pairs at $v = 0.5$ and at $v = 2$.
\end{problem}

\begin{solution}
The condition for $(x_a, t_a) \preceqLR (x_b, t_b)$ is
$t_a < t_b$ and $\abs{x_a - x_b} \le v\,(t_b - t_a)$.  Time
differences in the snapshot set come from $\Delta t \in \{0.5, 1.0, 1.5\}$
with multiplicities $3$, $2$, $1$ respectively (a standard count of
ordered pairs $\binom{4}{2} = 6$).  Spatial gaps on a five-site chain
are $0, 1, 2, 3, 4$, with multiplicities $5, 4, 3, 2, 1$.

\textbf{(a) $v = 0.5$.}  At $\Delta t = 0.5, 1.0, 1.5$ the maximum
admissible spatial gap is $0.25, 0.5, 0.75$ respectively --- below
$1$ in every case.  Only same-site (gap $= 0$) pairs are LR-related:
\[
  \#\{\text{pairs}\} \;=\; (3 + 2 + 1) \times 5 \;=\; 30.
\]

\textbf{(b) $v = 2$.}  At $\Delta t = 0.5$, gap $\le 1$ admits
$5 + 4 = 9$ spatial pairs; at $\Delta t = 1.0$, gap $\le 2$ admits
$5 + 4 + 3 = 12$; at $\Delta t = 1.5$, gap $\le 3$ admits
$5 + 4 + 3 + 2 = 14$.  Total:
\[
  3 \cdot 9 \;+\; 2 \cdot 12 \;+\; 1 \cdot 14
  \;=\; 27 + 24 + 14 \;=\; 65.
\]
A wider LR cone admits more than twice as many pairs as the narrower
one.
\end{solution}

\begin{problem}
Show $\preceqLR \subseteq \preceqcs$ on the regular chain, hence
$\tau(\LR, \cs) = +1$.
\end{problem}

\begin{solution}
On the regular $1{+}1$D chain $\preceqcs$ is the time-only order:
$(A, s) \preceqcs (B, t)$ iff $s < t$.  Now suppose
$(A, s) \preceqLR (B, t)$.  By the definition of $\preceqLR$ this
requires $s < t$ (and a spatial-gap condition that we do not need
here).  Hence $(A, s) \preceqcs (B, t)$.  This gives the inclusion
$\preceqLR \subseteq \preceqcs$ as relations on the same label set.

Consequently every pair comparable in both orders relates the same
way (because $\preceqLR$ is just a sub-relation of $\preceqcs$), so
$n_{cc} = \abs{\preceqLR}$, $n_{dd} = 0$, and
$\tau(\LR, \cs) = (\abs{\preceqLR} - 0)/(\abs{\preceqLR} + 0) = +1$.
This is the trivial-agreement caveat: the Kendall coefficient is
forced to its extremal value by the lattice geometry, not by any
property of the state.
\end{solution}

% =====================================================================
\section{Comparison statistics}

\begin{problem}
Compute $n_{cc}, n_{dd}, n_a^o, n_b^o, \tau$ for the two specified
partial orders on $\{1, 2, 3, 4, 5\}$.
\end{problem}

\begin{solution}
Transitive closures:
\begin{itemize}[nosep]
\item $\preceq_a$: $1 \prec 2$, $1 \prec 3$, $2 \prec 3$, $4 \prec 5$.
\item $\preceq_b$: $1 \prec 3$, $2 \prec 4$, $2 \prec 5$, $4 \prec 5$.
\end{itemize}
The ten unordered pairs:
\begin{center}
\renewcommand{\arraystretch}{1.05}
\begin{tabular}{c|c|c|l}
pair & in $\preceq_a$ & in $\preceq_b$ & class \\\hline
$\{1, 2\}$ & $1 \prec_a 2$  & incomparable    & $\preceq_a$ only \\
$\{1, 3\}$ & $1 \prec_a 3$  & $1 \prec_b 3$    & concordant \\
$\{1, 4\}$ & incomparable    & incomparable    & both incomparable \\
$\{1, 5\}$ & incomparable    & incomparable    & both incomparable \\
$\{2, 3\}$ & $2 \prec_a 3$  & incomparable    & $\preceq_a$ only \\
$\{2, 4\}$ & incomparable    & $2 \prec_b 4$    & $\preceq_b$ only \\
$\{2, 5\}$ & incomparable    & $2 \prec_b 5$    & $\preceq_b$ only \\
$\{3, 4\}$ & incomparable    & incomparable    & both incomparable \\
$\{3, 5\}$ & incomparable    & incomparable    & both incomparable \\
$\{4, 5\}$ & $4 \prec_a 5$  & $4 \prec_b 5$    & concordant \\
\end{tabular}
\end{center}
Counts: $n_{cc} = 2$, $n_{dd} = 0$, $n_a^o = 2$, $n_b^o = 2$.
Therefore
\[
  \tau_{ab} \;=\; \frac{2 - 0}{2 + 0} \;=\; +1.
\]
Note the contrast with raw pair counts: only $4$ of the $10$ pairs
are comparable in both orders, but on those $4$ they agree
unanimously.  This illustrates why the asymmetric counts $n_a^o$ and
$n_b^o$ carry information beyond $\tau$: they tell you \emph{how
much overlap} there is to compare in the first place.
\end{solution}

% =====================================================================
\section{Interpretation}

\begin{problem}
State the difference between the three falsification criteria and
give one example of each.
\end{problem}

\begin{solution}
\textbf{Strong falsification (criterion 1).}  $\preceqmaj$ relates a
pair whose endpoints lie outside the LR cone at the corresponding
time difference.  Concretely: there exist labels $(A, s), (B, t)$
with $\bm\lambda_A(s) \preceq \bm\lambda_B(t)$ but
$\mathrm{gap}(A, B) > v_{\LR}\, (t - s)$.  Such a pair would require
a faster-than-Lieb--Robinson information channel to be consistent
with the hypothesis, so it directly refutes (H).

\textit{Example.}  In Phase~5 (Table~2 of the paper), the count
$n_{\maj \notin \LR}$ at $v = 1$ is positive in every regime; each
contributing pair is a strong-falsification observation.

\smallskip
\textbf{Weak falsification (criterion 2).}  $\preceqmaj$ disagrees
with a non-trivially-defined $\preceqcs$ on a non-vanishing fraction
of comparable pairs.  This is meaningful only when $\preceqcs$ has
within-time-slice structure --- i.e., on a non-regular causet --- and
asks whether the prescribed causet structure controls the
entanglement order.

\textit{Example.}  On a $1{+}1$D causal dynamical triangulation
with multi-vertex antichains, find a pair related by $\preceqcs$
but not $\preceqmaj$, where the cs relation is genuinely
within-slice (not implied by the time order).  This is the
Phase~6 test of the implementation plan.

\smallskip
\textbf{Trivial agreement (criterion 3).}  $\preceqcs$ is forced to
coincide with $\preceqLR$ by lattice regularity.  Any agreement
between $\preceqmaj$ and $\preceqcs$ is then uninformative because
$\preceqcs$ has no extra structure to confirm against.

\textit{Example.}  On the regular chain, $\preceqcs$ reduces to the
time-only order and contains $\preceqLR$ as a strict subset.  An
observation $\tau(\maj, \cs) \approx +1$ in this geometry would
\emph{not} support (H), because $\preceqcs$ is here a coarse
total-time order that any half-decent state-derived ordering will
mostly agree with.

\smallskip
\textit{Strong but not weak.}  A pair on the regular chain whose
spatial gap exceeds $v_{\LR}(t - s)$ and is $\preceqmaj$-related is a
strong-falsification observation.  It is not a weak-falsification
observation because the regular chain has trivial $\preceqcs$
structure --- criterion 2 is not even posable in this geometry.
\end{solution}

\begin{problem}[$\star$]
Why is the trend $\tau(\maj, \LR)\!\downarrow$ as $v\!\uparrow$
counterintuitive, and what does it imply?
\end{problem}

\begin{solution}\hfill
\begin{enumerate}[label=(\alph*),leftmargin=*]
\item \emph{Naive expectation.}  Widening $v$ admits more pairs into
      $\preceqLR$.  If those new pairs were drawn at random from the
      $\preceqmaj$-comparable set with the same agreement rate as the
      existing pairs, the both-comparable count $n_{cc} + n_{dd}$
      would grow and $\tau$ would stay roughly constant or, if the
      new pairs were even better aligned with $\preceqmaj$ than the
      old ones, rise.  The intuition that ``a bigger
      cone~$=$~more~causal~pairs~$=$~closer~to~capturing~$\preceqmaj$''
      predicts $\tau\!\uparrow$.
\item \emph{What we see.}  $\tau$ drifts \emph{down}.  The pairs
      newly admitted by widening the cone do not merely fail to
      improve the agreement; they actively pull $\tau$ down by being
      net discordant with $\preceqmaj$ on the both-comparable subset
      (recall pairs that are LR-comparable but not maj-comparable
      sit in $n_b^o$ and do not affect $\tau$ directly).  So at
      least some of the newly-admitted ``causal'' pairs are
      ordered by $\preceqmaj$ in the opposite direction from the
      time order.
\item \emph{Picture.}  Imagine $\preceqmaj$ as carrying its own
      effective ``cone'' of slope $v_{\maj}$.  Pairs strictly inside
      that cone tend to be concordant with the time direction.
      Pairs in the LR ring outside $v_{\maj}$ but inside $v$
      contribute mostly discordant relations.  As $v$ grows past
      $v_{\maj}$, the discordant ring grows faster than the
      concordant interior, so $\tau$ falls.
\end{enumerate}
\end{solution}

\begin{problem}[$\star\star$]
Suggest two refinements of (H) that could survive the Phase~5
result.
\end{problem}

\begin{solution}
There is no unique correct answer.  Two natural refinements:

\textbf{Refinement~1: bounded-length cuts.}
Restrict the cut family $\mathcal{F}$ to intervals of length
$\le L_{\max}$.  Hypothesis (H${}_1$): the three orders coincide
when restricted to short-interval cuts.  Test: re-run the agreement
scan at varying $L_{\max}$.  Support: a sharp rise in
$\tau(\maj, \LR)$ as $L_{\max}$ shrinks would suggest the
discordance is concentrated in long-interval cuts, perhaps because
long intervals see boundary effects or accumulate local truncation
errors.

\textbf{Refinement~2: charge-sectorized majorization.}
Compute Schmidt spectra restricted to fixed total-charge sectors of
each bipartition --- i.e., for each region $A$ and each value $q$
of $\sum_{n \in A} \tfrac{1 - Z_n}{2}$, build the Schmidt spectrum
$\bm\lambda_A^{(q)}(t)$ of the projection onto the $q$-sector.
Define $\preceqmaj^{(q)}$ accordingly.  Hypothesis (H${}_2$): the
sectorized maj order $\preceqmaj^{(q)}$ coincides with $\preceqLR$
on each sector.  Test: re-run with quantum-number-resolved spectra,
already supported by the underlying ITensor implementation.
Support: the global mixing of charge sectors might be the source of
the $\sim 50\%$ overflow; resolving it could push the overflow
fraction to small.

Other plausible refinements include: time-coarse-graining (group
consecutive snapshots), substituting Hasse-graph edit distance for
exact-equality of orders, and weighted Kendall coefficients that
discount pairs near the cone boundary.
\end{solution}

% =====================================================================
\section{Computational exercises (optional)}

\begin{problem}
Run \texttt{lightcone\_vs\_majorization.py} and verify Table~1.
\end{problem}

\begin{solution}
The expected wall-clock time on a single laptop CPU is roughly $80$
seconds.  Output to within $\sim 10^{-4}$ of the paper's Table~1 is
expected; small discrepancies arise from compiler-level
numerical differences (LAPACK / BLAS implementation, FMA contraction)
and Krylov-subspace basis ordering and are not physically
significant.  A run that disagrees by more than $\sim 10^{-3}$ on any
$\tau$ value indicates a configuration mismatch (commonly: bond
dimension or DMRG sweep count) and warrants investigation.
\end{solution}

\begin{problem}[$\star$]
Add an intermediate $m/g = 1.0$ regime and compare.
\end{problem}

\begin{solution}
Adding a regime is a one-line modification to \texttt{lightcone\_vs\_majorization.py}:
add a fourth call to \texttt{scan\_vlr} with \texttt{m\_over\_g=1.0}
between regimes A and B.

The expected qualitative finding is monotone behavior: the
intermediate $m/g = 1.0$ produces $\tau(\maj, \LR)$ values
\emph{between} the light-quark ($m/g = 0.5$) and heavy-quark
($m/g = 5.0$) values at every $v$, with $\tau$ rising smoothly with
$m/g$.  The physical intuition is that increasing $m/g$ stiffens
the flux tube and suppresses pair production; the entanglement
spreading slows; the Schmidt spectra evolve more in time than across
cuts; $\preceqmaj$ tracks the time order more closely; and
$\tau(\maj, \LR)$ rises.

A non-monotone observation would be surprising and physically
interesting --- it would imply a non-trivial competition between
two effects (e.g.\ string fluctuations vs.\ pair production) that
maximally disorganizes the spectra at intermediate $m/g$.
\end{solution}

\end{document}
